\section{A catalogue of extremely amenable groups}\label{sec:ea-catalogue}

Extremely amenable groups arise from concentration of measure, from structural
Ramsey theory, and as stabilisers inside groups whose universal minimal flow
is metrizable. This section collects the main examples; groups and
topologies are as in the cited sources.

\begin{theorem}[Pestov \cite{Pestov2002}]\label{thm:pestov-urysohn}
\uses{def:ea}
The isometry group $\mathrm{Iso}(\mathbb{U})$ of the Urysohn universal
complete separable metric space $\mathbb{U}$, with the compact-open topology, is
extremely amenable. The same holds for sufficiently homogeneous generalised
Urysohn metric spaces, and together with a theorem of Uspenskij this shows that
every topological group embeds into an extremely amenable group.
\end{theorem}

\begin{theorem}[Giordano--Pestov \cite{GiordanoPestov2007}]\label{thm:giordano-pestov}
\uses{def:ea}
\begin{enumerate}
\item The group of measure-preserving automorphisms of a standard Lebesgue
  space with a non-atomic measure, with the weak topology, is extremely
  amenable.
\item The isometry groups of the spaces $L^p(0,1)$, $1\leq p<\infty$, are
  extremely amenable. This extends the case $p=2$ of
  Theorem~\ref{thm:gromov-milman}.
\item A measure class preserving equivalence relation $R$ on a standard Borel
  space is amenable if and only if its full group $[R]$, with the uniform
  topology, is extremely amenable.
\item A von Neumann algebra is approximately finite-dimensional if and only if
  its unitary group, with the strong topology, is the product of an extremely
  amenable group and a compact group.
\end{enumerate}
\end{theorem}

\begin{remark}[Groups of measurable maps]\label{rem:l0}
\uses{def:ea, thm:glasner-monothetic}
The group $L^0(I,U(1))$ of measurable maps from the unit interval to the circle
group, with the topology of convergence in measure, is extremely amenable. This
was proved by Glasner and, independently and unpublished, by Furstenberg and
Weiss; see \cite{GiordanoPestov2007}. Glasner also constructed an
infinite monothetic extremely amenable Polish group
(Theorem~\ref{thm:glasner-monothetic}).
\end{remark}

\begin{remark}[Summary of the examples in this encyclopaedia]\label{rem:ea-summary}
\uses{rem:herer-christensen, thm:gromov-milman, rem:levy, thm:pestov-ea, thm:kpt, ex:kpt-computations, ex:standard-chain, rem:lelek-method, cor:lc-not-ea}
\begin{enumerate}
\item The groups of Herer and Christensen (Remark~\ref{rem:herer-christensen}),
  historically the first examples.
\item The unitary group of $\ell^2$ with the strong operator topology
  (Theorem~\ref{thm:gromov-milman}), and more generally Lévy groups
  (Remark~\ref{rem:levy}).
\item $\mathrm{Aut}(\mathbb{Q},<)$ and $\mathrm{Homeo}_+[0,1]$
  (Theorem~\ref{thm:pestov-ea}).
\item Automorphism groups of Fraïssé limits of Fraïssé order classes with the
  Ramsey property (Theorem~\ref{thm:kpt}), such as the order expansions in
  Example~\ref{ex:kpt-computations}.
\item The stabiliser of the standard chain in $\mathrm{Homeo}(E)$, for $E$ the
  Cantor set (Example~\ref{ex:standard-chain}), and the stabiliser of a
  generic connected chain in the homeomorphism group of the Lelek fan
  (Remark~\ref{rem:lelek-method}).
\end{enumerate}
By Corollary~\ref{cor:lc-not-ea}, none of these groups is locally compact;
according to \cite{GiordanoPestov2007}, this was first shown by Granirer and
Lau.
\end{remark}
